\documentclass[11pt,a4paper]{article}

\usepackage[utf8]{inputenc}
\usepackage[T1]{fontenc}
\usepackage[french]{babel}
\usepackage{geometry}
\usepackage[hidelinks]{hyperref}
\usepackage{booktabs}
\usepackage{enumitem}
\usepackage{titlesec}
\usepackage{fancyhdr}
\usepackage{parskip}
\usepackage{array}
\usepackage{tabularx}
\usepackage{lmodern}

\geometry{top=2.5cm, bottom=2.5cm, left=2.8cm, right=2.8cm}

\pagestyle{fancy}
\fancyhf{}
\fancyhead[L]{\small\textit{git-archaeologist -- Rapport DevOps \& AWS}}
\fancyhead[R]{\small\textit{Maria Khvatova}}
\fancyfoot[C]{\thepage}
\renewcommand{\headrulewidth}{0.4pt}

\titleformat{\section}{\large\bfseries}{\thesection}{1em}{}[\vspace{-4pt}\rule{\linewidth}{0.4pt}]
\titleformat{\subsection}{\normalsize\bfseries}{\thesubsection}{1em}{}

\setlength{\parskip}{6pt}
\setlength{\parindent}{0pt}

\newcommand{\code}[1]{\texttt{#1}}

\begin{document}

% ── Title page ────────────────────────────────────────────────────────────────
\begin{titlepage}
  \centering
  \vspace*{3cm}
  {\Huge\bfseries git-archaeologist\par}
  \vspace{0.5cm}
  {\large Service cloud d'analyse de l'historique Git par LLM\par}
  \vspace{2cm}
  \rule{0.5\linewidth}{0.5pt}\par
  \vspace{1.5cm}
  {\large Rapport de projet -- DevOps \& AWS\par}
  \vspace{0.5cm}
  {\normalsize EPITA ING2 -- 2025/2026\par}
  \vspace{1cm}
  {\normalsize Maria Khvatova\par}
  \vspace{0.3cm}
  {\small \url{https://github.com/khvatovat/git-archaeologist}\par}
  \vfill
  {\small Juin 2026}
\end{titlepage}

\tableofcontents
\newpage

% ─────────────────────────────────────────────────────────────────────────────
\section{Presentation de l'application}

\subsection{Contexte}

\code{git-archaeologist} est un outil d'analyse de l'historique Git d'un fichier dans un depot GitHub. Pour chaque fichier analyse, il recupere les commits et pull requests associes, les regroupe par periode (appel\'{e}es \textit{eres}), g\'{e}n\`{e}re un r\'{e}sum\'{e} narratif de chaque \`{e}re avec Claude Haiku, puis synth\'{e}tise l'ensemble avec Claude Sonnet pour produire un rapport lisible sur l'\'{e}volution du fichier.

Ce projet transforme cet outil CLI en service cloud h\'{e}berg\'{e} sur AWS~:

\begin{itemize}[noitemsep]
  \item L'utilisateur enregistre un d\'{e}pot GitHub via une interface web statique (S3).
  \item GitHub envoie un webhook \`{a} chaque push~; l'API valide la signature HMAC-SHA256, cr\'{e}e une entr\'{e}e dans DynamoDB et pousse un message dans SQS.
  \item Un worker consomme la file SQS, ex\'{e}cute le pipeline d'analyse et \'{e}crit le r\'{e}sultat dans DynamoDB.
  \item L'interface web permet de consulter les analyses par d\'{e}pot.
\end{itemize}

La logique d'analyse IA est r\'{e}utilis\'{e}e telle quelle depuis l'outil CLI. L'objectif du projet est l'infrastructure, l'automatisation, la s\'{e}curit\'{e} et l'observabilit\'{e}.

\subsection{Stack applicative}

\begin{tabularx}{\linewidth}{lX}
\toprule
Composant & Technologies \\
\midrule
API & FastAPI (Python 3.11), Uvicorn, Pydantic \\
Worker & Python 3.11, boto3, Anthropic SDK \\
Frontend & HTML/JS statique heberg\'{e} sur S3 \\
LLM & Claude Haiku (r\'{e}sum\'{e}s d'eres), Claude Sonnet (synth\`{e}se) \\
Tests & pytest, moto (mock AWS), coverage \\
\bottomrule
\end{tabularx}

% ─────────────────────────────────────────────────────────────────────────────
\section{Architecture AWS}

\subsection{Vue d'ensemble}

\begin{verbatim}
  Internet
     |
  ALB (public, port 80)              us-east-1
     |
  ECS Fargate -- api           (subnet prive)
     |-- POST /repos            --> DynamoDB repos
     |-- GET  /repos            <-- DynamoDB repos
     |-- POST /webhook          --> DynamoDB jobs + SQS
     |-- GET  /repos/{id}/history  <-- DynamoDB jobs (GSI)
     |
  SQS --> ECS Fargate -- worker  (subnet prive)
               |
          DynamoDB jobs   (statut, resultat)
               |
          Secrets Manager --> ANTHROPIC_API_KEY
                          --> GITHUB_TOKEN
\end{verbatim}

\subsection{Reseau (module \code{network})}

Le VPC couvre deux zones de disponibilit\'{e}. Chaque AZ contient un subnet public (ALB, NAT Gateway) et un subnet priv\'{e} (t\^{a}ches ECS). Les subnets priv\'{e}s n'ont pas de route vers l'Internet Gateway~; le trafic sortant passe par la NAT Gateway. Une seule NAT Gateway est provisionn\'{e}e (optimisation de co\^{u}t acceptable en d\'{e}veloppement~; en production, une par AZ garantit la r\'{e}silience).

\subsection{S\'{e}curit\'{e} r\'{e}seau (module \code{security})}

Trois security groups strictement d\'{e}limit\'{e}s~:

\begin{tabularx}{\linewidth}{lXX}
\toprule
Security group & Inbound & Outbound \\
\midrule
\code{alb-sg} & \code{0.0.0.0/0 :80} & Tout \\
\code{api-sg} & \code{alb-sg :8080} uniquement & Tout (APIs AWS via NAT) \\
\code{worker-sg} & Aucun & Tout \\
\bottomrule
\end{tabularx}

L'API n'est jamais accessible directement depuis internet, uniquement via l'ALB. Le worker n'a aucun port ouvert en entr\'{e}e.

\subsection{Pourquoi DynamoDB et SQS}

\textbf{DynamoDB} est retenu car tous les acc\`{e}s sont par cl\'{e}~: recherche d'un job par \code{job\_id}, liste des jobs d'un d\'{e}pot via GSI sur \code{repo\_id+created\_at}, enregistrement d'un d\'{e}pot par \code{repo\_id}. Il n'y a pas de jointures ni d'agr\'{e}gations. Le mode \code{PAY\_PER\_REQUEST} rend le co\^{u}t n\'{e}gligeable en d\'{e}veloppement. RDS apporterait une complexit\'{e} op\'{e}rationnelle (backups, patches, connection pooling) sans aucun b\'{e}n\'{e}fice pour ces patterns d'acc\`{e}s.

\textbf{SQS} est retenu car les jobs d'analyse sont longs (de l'ordre de la minute) et naturellement asynchrones. Le pattern producteur/consommateur est exactement le cas d'usage de SQS. Le \textit{visibility timeout} et la dead-letter queue (3 tentatives) g\`{e}rent les \'{e}checs sans perte de donn\'{e}es. Step Functions serait disproportionn\'{e} pour un flux \`{a} une seule \'{e}tape.

\subsection{IAM -- principe du moindre privil\`{e}ge}

Chaque service n'a acc\`{e}s qu'aux ressources dont il a strictement besoin~:

\begin{tabularx}{\linewidth}{lX}
\toprule
R\^{o}le & Actions autoris\'{e}es \\
\midrule
\code{api-task-role} & DynamoDB \code{PutItem/GetItem/Scan/Query} (repos + jobs), SQS \code{SendMessage} \\
\code{worker-task-role} & DynamoDB \code{GetItem/UpdateItem} (jobs uniquement), SQS \code{ReceiveMessage/Delete\-Message}, Secrets Manager \code{GetSecretValue} \\
\code{execution-role} & ECR pull, CloudWatch Logs write, Secrets Manager \code{GetSecretValue} (injection dans le container) \\
\bottomrule
\end{tabularx}

Le worker ne peut ni lire ni \'{e}crire la table des d\'{e}pots. L'API ne peut ni lire les secrets ni acc\'{e}der aux champs de mise \`{a} jour du worker. Les ARN des ressources sont explicitement sp\'{e}cifi\'{e}s dans les politiques (pas de wildcard \code{*} sur \code{Resource}).

\subsection{Secrets Manager}

La cl\'{e} Anthropic et le token GitHub sont stock\'{e}s dans AWS Secrets Manager. Ils sont inject\'{e}s dans le container worker via le champ \code{secrets} de la d\'{e}finition de t\^{a}che ECS, pas comme variables d'environnement d\'{e}finies dans Terraform. Les valeurs n'apparaissent donc ni dans l'\'{e}tat Terraform, ni dans les logs, ni dans le code source.

\subsection{V\'{e}rification de signature webhook}

Chaque webhook GitHub entrant est v\'{e}rifi\'{e} par HMAC-SHA256 avec le secret par d\'{e}pot stock\'{e} dans DynamoDB. La comparaison utilise \code{hmac.compare\_digest} pour pr\'{e}venir les attaques temporelles. Les requ\^{e}tes qui \'{e}chouent la v\'{e}rification retournent HTTP 401 avant toute \'{e}criture DynamoDB.

% ─────────────────────────────────────────────────────────────────────────────
\section{Infrastructure as Code -- Terraform}

\subsection{Structure modulaire}

L'infrastructure est organis\'{e}e en 7 modules ind\'{e}pendants, c\^{a}bl\'{e}s dans \code{terraform/envs/dev/main.tf}~:

\begin{tabularx}{\linewidth}{lX}
\toprule
Module & Ressources cr\'{e}\'{e}es \\
\midrule
\code{network} & VPC, 2 subnets publics, 2 subnets priv\'{e}s, IGW, NAT Gateway, tables de routage \\
\code{security} & Security groups (ALB, API, worker), Secrets Manager (cl\'{e} Anthropic, token GitHub) \\
\code{messaging} & File SQS + dead-letter queue (redrive apr\`{e}s 3 tentatives) \\
\code{data} & Tables DynamoDB \code{repos} et \code{jobs} avec chiffrement SSE activ\'{e}, index secondaires globaux \\
\code{ecs} & Cluster ECS Fargate, d\'{e}p\^{o}ts ECR (tags immuables), groupes de logs CloudWatch, r\^{o}les IAM, ALB, d\'{e}finitions de t\^{a}ches et services \\
\code{frontend} & Bucket S3 (site statique), politique publique en lecture, upload des fichiers HTML \\
\code{observability} & Dashboard CloudWatch (5 widgets), 2 alarmes \\
\bottomrule
\end{tabularx}

\subsection{Backend d'\'{e}tat}

L'\'{e}tat Terraform est stock\'{e} dans un bucket S3 avec versioning activ\'{e}. Le verrouillage est assur\'{e} par une table DynamoDB (\code{terraform-locks}), ce qui permet l'ex\'{e}cution du pipeline CI/CD sans risque de corruption de l'\'{e}tat.

\subsection{Choix ECS Fargate}

ECS Fargate a \'{e}t\'{e} retenu plut\^{o}t qu'EC2 ou Lambda~:

\begin{itemize}[noitemsep]
  \item \textbf{EC2 + Auto Scaling}~: n\'{e}cessite de g\'{e}rer les AMI, patches OS et placement d'instances. La charge ne justifie pas cette complexit\'{e}.
  \item \textbf{Lambda}~: la limite d'ex\'{e}cution \`{a} 15 minutes est incompatible avec les jobs d'analyse qui peuvent durer plusieurs minutes selon la taille du d\'{e}pot. La gestion des d\'{e}pendances (Anthropic SDK, boto3) via Lambda layers serait aussi plus lourde.
  \item \textbf{ECS Fargate}~: pas de gestion de serveurs, facturation \`{a} la seconde, isolation native des containers, int\'{e}gration directe avec CloudWatch Logs, Secrets Manager et SQS.
\end{itemize}

% ─────────────────────────────────────────────────────────────────────────────
\section{CI/CD -- GitHub Actions}

Trois workflows ind\'{e}pendants sont mis en place, chacun avec un p\'{e}rim\`{e}tre diff\'{e}rent.

\subsection{Pipeline A -- Application (\code{app.yml})}

D\'{e}clench\'{e} sur push ou pull request vers \code{main}, uniquement si des fichiers sous \code{api/}, \code{worker/} ou des \code{*.py} \`{a} la racine ont chang\'{e}.

\begin{enumerate}[noitemsep]
  \item \textbf{Unit Tests}~: \code{pytest api/tests/ worker/tests/} avec \code{coverage}. Les d\'{e}pendances AWS sont mock\'{e}es avec \code{moto} (DynamoDB, SQS, Secrets Manager).
  \item \textbf{SonarCloud}~: analyse de qualit\'{e} de code avec upload du rapport de couverture (push vers \code{main} uniquement).
  \item \textbf{Trivy Filesystem Scan}~: scan des vuln\'{e}rabilit\'{e}s CRITICAL et HIGH sur le syst\`{e}me de fichiers, r\'{e}sultats en SARIF.
  \item \textbf{Build, Scan \& Deploy} (push vers \code{main} uniquement)~: build Docker API, scan Trivy de l'image, push ECR~; idem pour le worker~; puis \code{ecs update-service --force-new-deployment} sur les deux services.
\end{enumerate}

\subsection{Pipeline B -- Infrastructure (\code{terraform.yml})}

D\'{e}clench\'{e} uniquement si des fichiers sous \code{terraform/} ont chang\'{e}.

\begin{enumerate}[noitemsep]
  \item \textbf{Format \& Validate}~: \code{terraform fmt -check} et \code{terraform validate}.
  \item \textbf{Checkov}~: scan de s\'{e}curit\'{e} IaC (r\`{e}gles Terraform), r\'{e}sultats en SARIF.
  \item \textbf{Terraform Plan} (pull requests uniquement)~: g\'{e}n\`{e}re le plan et le publie en commentaire sur la PR pour revue manuelle avant merge.
  \item \textbf{Terraform Apply} (push vers \code{main})~: \code{apply -auto-approve} avec variables sensibles inject\'{e}es depuis les secrets GitHub.
  \item \textbf{Ansible Post-Deploy}~: apr\`{e}s l'apply, un playbook attend que le health check de l'API r\'{e}ponde, synchronise les assets frontend vers S3, et soumet un job de test pour valider la cha\^{i}ne compl\`{e}te.
\end{enumerate}

\subsection{Pipeline C -- S\'{e}curit\'{e} (\code{security.yml})}

D\'{e}clench\'{e} sur chaque push vers \code{main}, chaque PR, et hebdomadairement (lundi 6h UTC).

\begin{enumerate}[noitemsep]
  \item \textbf{Semgrep SAST}~: analyse statique avec les r\`{e}gles \code{p/python} et \code{p/owasp-top-ten}. Couvre Python, Terraform, Dockerfile et YAML. Bloquant sur les findings critiques.
  \item \textbf{pip-audit}~: scan des d\'{e}pendances Python contre la base CVE, s\'{e}par\'{e}ment pour l'API et le worker.
  \item \textbf{TruffleHog}~: scan de l'historique Git complet pour d\'{e}tecter des secrets committ\'{e}s par erreur. Mode \code{--only-verified} pour limiter les faux positifs.
\end{enumerate}

\subsection{Couverture des outils de s\'{e}curit\'{e}}

\begin{tabularx}{\linewidth}{llX}
\toprule
Outil & Pipeline & P\'{e}rim\`{e}tre \\
\midrule
Semgrep & C & SAST Python + IaC, r\`{e}gles OWASP Top 10 \\
Checkov & B & Mauvaises configurations Terraform (chiffrement, IAM, r\'{e}seau) \\
Trivy & A & CVE sur filesystem et images Docker \\
TruffleHog & C & D\'{e}tection de secrets dans l'historique Git \\
pip-audit & C & CVE dans les d\'{e}pendances Python \\
SonarCloud & A & Qualit\'{e} de code, dette technique, couverture de tests \\
\bottomrule
\end{tabularx}

% ─────────────────────────────────────────────────────────────────────────────
\section{Automatisation -- Ansible}

Ansible est utilis\'{e} en post-d\'{e}ploiement pour valider que l'infrastructure est op\'{e}rationnelle et synchroniser les assets statiques. Le playbook (\code{ansible/deploy.yml}) s'ex\'{e}cute sur le runner GitHub Actions apr\`{e}s \code{terraform apply}~:

\begin{enumerate}[noitemsep]
  \item R\'{e}cup\`{e}re le DNS de l'ALB via \code{terraform output}.
  \item Attend le passage du health check \code{GET /health} avec retry (10 tentatives, 15 secondes d'intervalle).
  \item Synchronise les fichiers HTML du frontend vers le bucket S3.
  \item Soumet un job de test via \code{POST /jobs} pour valider la cha\^{i}ne API$\to$SQS$\to$worker.
\end{enumerate}

Ce choix \'{e}vite de dupliquer la logique de d\'{e}ploiement des assets dans Terraform et fournit une validation fonctionnelle automatis\'{e}e apr\`{e}s chaque d\'{e}ploiement.

% ─────────────────────────────────────────────────────────────────────────────
\section{Observabilit\'{e}}

\subsection{Logs}

Les logs de l'API et du worker sont stream\'{e}s vers CloudWatch Logs via le driver \code{awslogs} configur\'{e} dans les d\'{e}finitions de t\^{a}ches ECS (\code{/ecs/git-archaeologist-dev/api} et \code{/ecs/git-archaeologist-dev/worker}). R\'{e}tention fix\'{e}e \`{a} 14 jours.

\subsection{Dashboard CloudWatch}

Un dashboard centralise cinq m\'{e}triques cl\'{e}s~:

\begin{tabularx}{\linewidth}{lX}
\toprule
Widget & M\'{e}trique \\
\midrule
ALB Request Count & Nombre de requ\^{e}tes par minute (SUM) \\
ALB Response Time p95 & Latence au 95\textsuperscript{e} percentile \\
SQS Queue Depth & Messages visibles dans la file (backlog de jobs) \\
Worker Task Count & T\^{a}ches ECS en cours d'ex\'{e}cution (Container Insights) \\
ALB 5xx Error Count & Erreurs serveur par minute (SUM) \\
\bottomrule
\end{tabularx}

\subsection{Alarmes}

\begin{itemize}[noitemsep]
  \item \textbf{queue-depth-high}~: file SQS $>$ 20 messages pendant 2 minutes cons\'{e}cutives -- le worker est bloqu\'{e} ou sous-dimensionn\'{e}.
  \item \textbf{api-5xx-high}~: plus de 10 erreurs 5xx par minute pendant 2 minutes -- probl\`{e}me applicatif ou infrastructure.
\end{itemize}

% ─────────────────────────────────────────────────────────────────────────────
\section{Limites et am\'{e}liorations possibles}

\begin{tabularx}{\linewidth}{lX}
\toprule
Limite actuelle & Mitigation en production \\
\midrule
NAT Gateway unique & Une NAT GW par AZ pour \'{e}viter le point de d\'{e}faillance sur la sortie priv\'{e}e \\
ALB en HTTP uniquement & Certificat ACM + listener HTTPS + redirection HTTP vers HTTPS \\
Worker non scalable & ECS Application Auto Scaling sur la m\'{e}trique de profondeur de file SQS \\
Credentials AWS temporaires & OIDC federation GitHub$\to$IAM (pas de secrets statiques \`{a} stocker) \\
Pas d'auth sur l'API & API Gateway + Cognito ou cl\'{e} API ; le secret webhook est la seule protection actuelle \\
\bottomrule
\end{tabularx}

% ─────────────────────────────────────────────────────────────────────────────
\section{Estimation des co\^{u}ts (environnement dev)}

\begin{tabularx}{\linewidth}{lX}
\toprule
Ressource & Co\^{u}t mensuel estim\'{e} \\
\midrule
ECS Fargate API (0,25 vCPU / 0,5~GB, 1 t\^{a}che) & \$6 \\
ECS Fargate Worker (0,5 vCPU / 1~GB, 1 t\^{a}che) & \$11 \\
NAT Gateway (1, trafic minimal) & \$35 \\
ALB & \$16 \\
Secrets Manager (2 secrets) & \$0,80 \\
CloudWatch (logs + dashboard) & \$3 \\
DynamoDB + SQS (trafic dev) & $<$ \$1 \\
\midrule
\textbf{Total} & \textbf{\textasciitilde\$73/mois} \\
\bottomrule
\end{tabularx}

La NAT Gateway repr\'{e}sente la majorit\'{e} du co\^{u}t. En production, des VPC Endpoints (S3, DynamoDB, SQS, Secrets Manager) r\'{e}duiraient le trafic qui la traverse.

\end{document}
